% Sample LaTeX Publication Template
% This template demonstrates the structure of a typical academic publication
% Compile with: pdflatex sample.tex && bibtex sample && pdflatex sample.tex && pdflatex sample.tex

\documentclass[12pt,a4paper]{article}

% Essential packages
\usepackage[utf8]{inputenc}
\usepackage[T1]{fontenc}
\usepackage{graphicx}           % For including figures
\usepackage{amsmath,amssymb}    % For mathematical equations
\usepackage{cite}               % For citation management
\usepackage{hyperref}           % For hyperlinks (should be loaded last)
\usepackage{booktabs}           % For professional-looking tables

% Page layout
\usepackage[margin=1in]{geometry}

% Title and author information
\title{Sample Publication Template: \\
       A Comprehensive Guide to Academic Writing}

\author{
    First Author\thanks{Corresponding author: first.author@university.edu} \\
    Department of Computer Science \\
    University Name \\
    \and
    Second Author \\
    Department of Engineering \\
    Another Institution
}

\date{\today}

\begin{document}

% Generate title page
\maketitle

% Abstract section - brief summary of the work
\begin{abstract}
This document serves as a sample template for creating academic publications using \LaTeX{}. It demonstrates proper document structure, including sections for abstract, introduction, methodology, results, conclusions, and references. The template also includes examples of common elements such as figures, tables, equations, and citations. This template is designed to be easily adaptable for various types of academic publications, including research papers, conference proceedings, and technical reports.

\noindent\textbf{Keywords:} LaTeX, Academic Writing, Publication Template, Document Formatting
\end{abstract}

% Start new page after abstract (optional)
\newpage

% Table of contents (optional, remove if not needed)
\tableofcontents
\newpage

% Main content sections
% Section 1: Introduction
\section{Introduction}
\label{sec:introduction}

The introduction provides background information and context for the research. It should clearly state the problem being addressed and explain why it is important. The introduction typically includes:

\begin{itemize}
    \item Background and motivation for the research
    \item Statement of the research problem or question
    \item Overview of the approach or methodology
    \item Summary of key contributions
    \item Organization of the paper
\end{itemize}

This sample document demonstrates the use of various \LaTeX{} features commonly used in academic publications. For example, we can cite previous work using BibTeX \cite{knuth1984texbook, lamport1994latex}. Cross-references to other sections (e.g., Section~\ref{sec:methods}) and figures (e.g., Figure~\ref{fig:sample}) are automatically numbered and updated.

% Section 2: Background/Related Work (optional)
\section{Related Work}
\label{sec:related}

This section discusses relevant prior research and situates your work within the existing literature. It demonstrates familiarity with the field and highlights how your work differs from or builds upon previous studies.

Recent advances in the field have been documented extensively \cite{goossens1994latex}. Multiple citations can be grouped together \cite{knuth1984texbook,lamport1994latex,goossens1994latex} for efficient referencing.

% Section 3: Methods/Methodology
\section{Methods}
\label{sec:methods}

The methods section describes the approach, techniques, and procedures used in the research. It should provide enough detail for others to reproduce your work.

\subsection{Experimental Setup}
\label{subsec:setup}

Describe the experimental setup, tools, and resources used. This subsection can include details about:
\begin{enumerate}
    \item Hardware and software specifications
    \item Data collection procedures
    \item Parameter settings and configurations
    \item Validation approaches
\end{enumerate}

\subsection{Mathematical Formulation}
\label{subsec:math}

Mathematical equations are a key component of technical publications. Here are examples of both inline and displayed equations.

Inline equations appear within the text, such as $E = mc^2$ or $f(x) = ax^2 + bx + c$.

Displayed equations are centered and numbered for reference:

\begin{equation}
\label{eq:quadratic}
x = \frac{-b \pm \sqrt{b^2 - 4ac}}{2a}
\end{equation}

Multi-line equations can be aligned at specific points:

\begin{align}
f(x) &= (x+1)^2 \label{eq:expand1} \\
     &= x^2 + 2x + 1 \label{eq:expand2}
\end{align}

We can reference Equation~\ref{eq:quadratic} or Equations~\ref{eq:expand1} and \ref{eq:expand2} in the text.

% Section 4: Results
\section{Results}
\label{sec:results}

The results section presents the findings of your research, typically using figures, tables, and descriptive text.

\subsection{Tables}
\label{subsec:tables}

Tables are useful for presenting structured data. Table~\ref{tab:results} shows an example of a professional-looking table using the \texttt{booktabs} package.

\begin{table}[ht]
\centering
\caption{Sample experimental results showing performance metrics}
\label{tab:results}
\begin{tabular}{lccc}
\toprule
Method & Accuracy (\%) & Precision (\%) & Recall (\%) \\
\midrule
Baseline & 75.3 & 72.1 & 78.5 \\
Proposed & 89.7 & 87.4 & 91.2 \\
State-of-art & 82.1 & 80.3 & 84.6 \\
\bottomrule
\end{tabular}
\end{table}

\subsection{Figures}
\label{subsec:figures}

Figures are essential for visualizing results. Figure~\ref{fig:sample} demonstrates how to include images in your document.

\begin{figure}[ht]
\centering
% Uncomment the line below when you have an actual figure file
% \includegraphics[width=0.6\textwidth]{figure.png}
% For now, we'll use a placeholder
\fbox{\parbox{0.6\textwidth}{\centering\vspace{2cm}
[Your Figure Here]\\
Replace with: \texttt{\textbackslash includegraphics[width=0.6\textbackslash textwidth]\{figure.png\}}
\vspace{2cm}}}
\caption{Sample figure demonstrating image placement. Replace this with your actual figure and caption.}
\label{fig:sample}
\end{figure}

% Section 5: Discussion (optional but recommended)
\section{Discussion}
\label{sec:discussion}

The discussion section interprets the results, explains their significance, and addresses limitations. It should:

\begin{itemize}
    \item Interpret the findings in context of the research question
    \item Compare results with previous work
    \item Discuss implications and applications
    \item Acknowledge limitations and potential sources of error
    \item Suggest directions for future research
\end{itemize}

% Section 6: Conclusion
\section{Conclusion}
\label{sec:conclusion}

The conclusion summarizes the key findings and contributions of the work. It should be concise and clearly state:

\begin{itemize}
    \item Main results and achievements
    \item Significance and impact of the work
    \item Final thoughts and future directions
\end{itemize}

This template provides a foundation for creating well-structured academic publications in \LaTeX{}. By following this structure and adapting it to your specific needs, you can produce professional-quality documents suitable for submission to journals, conferences, and other academic venues.

% Acknowledgments section (optional)
\section*{Acknowledgments}
We would like to thank [funding agency/institution] for supporting this research through [grant number]. We also acknowledge [collaborators/contributors] for their valuable input and assistance.

% Bibliography
% The bibliography is automatically generated from references.bib
% Compile sequence: pdflatex -> bibtex -> pdflatex -> pdflatex
\bibliographystyle{plain}  % Other styles: alpha, abbrv, ieeetr, acm, etc.
\bibliography{references}   % This refers to references.bib file

\end{document}
